\documentclass[11pt]{article}
\usepackage[a4paper,margin=24mm]{geometry}
\usepackage[T1]{fontenc}
\usepackage{amsmath,amssymb,amsthm,mathtools}
\usepackage{xurl}
\usepackage[hidelinks]{hyperref}
\allowdisplaybreaks
\setlength{\emergencystretch}{3em}
\newtheorem{theorem}{Theorem}
\newtheorem{lemma}[theorem]{Lemma}
\newcommand{\R}{\mathbb R}
\newcommand{\C}{\mathbb C}
\newcommand{\Pcal}{\mathcal P}
\newcommand{\norm}[1]{\left\lVert#1\right\rVert}
\title{A finite auxiliary-degree bound for the original native resolvent}
\author{}
\date{20 September 2026}
\begin{document}\maketitle

This gives an independent proof of equations (35)--(44) in the supplied
combined continuation. The arithmetic measure, the physical variable
$y$, its mass, the tilted parity measures and the original poles remain
unchanged. The exact finite native moments have already been identified
in RR1--16 \cite{RR}. Here the degree of the polynomial trial space is
bounded before running a convergence test. The full proof below supplies
the details of the Gamma kernel estimates and of the lower-envelope
comparison used by that degree bound. The actual lower and upper
envelopes are the previously proved programme results \cite{NP}; their
unevaluated arithmetic constants retain their definitions.
The Gauss/Radau comparison has its established human context in
Zimmerling, Druskin and Simoncini \cite{ZDS}, Sections 4--5.1.
Their positive finite block recurrence is not assumed for the original
arithmetic multiplication operator; NR6--7 give its direct native
minimum and exact error direction.

\section{The original arithmetic objects}

Keep the original stipulated quartet domain
\begin{equation}\tag{NR1}\label{nr:original}
\begin{gathered}
 \rho_*=\tfrac12+\delta+i\gamma,\qquad
 0<\delta<\tfrac12,\quad \gamma>2,\quad m_0\ge1,\quad g=2\xi,\\
 h(s)=\prod_{z\in\{\rho_*,\bar\rho_*,1-\rho_*,1-\bar\rho_*\}}(s-z)^{m_0},
 \qquad w_h(y)=\frac{|(g/h)(1/2+iy)|^2}{2\pi},\\
 m_k=w_h^{*k},\quad k=4\ell+1\ge9,\quad
 \sigma(y)=\frac{|\Gamma(1/4+iy/2)|^2}{2\pi}.
\end{gathered}
\end{equation}
No existence of an off-line zero is asserted by stipulating the original
quartet domain. The established original envelopes are
\begin{equation}\tag{NR2}\label{nr:envelope}
 \underline a_k(1+y^2)^{-M}\sigma(y)
 \le m_k(y)\le \overline a_k\sigma(y),\qquad M=21+4m_0,
\end{equation}
where, retaining the provider's constants,
\begin{equation}\tag{NR3}\label{nr:constants}
\begin{gathered}
 \alpha=\pi/2,\quad p_h=8m_0-9/2,\quad B_h=42+8m_0,\\
 c_\Gamma=\sqrt2e^{-7/6},\quad C_\Gamma=\sqrt{2\sqrt5}e^{2/3},\quad
 \vartheta_h=\int_{-1}^1w_h(y)\,dy,\quad
 M_h(\alpha)=\int_\R e^{\alpha y}w_h(y)\,dy,\\
 \underline a_k=
 \frac{c_h\vartheta_h^{k-3}e^{-\alpha(k-3)}(k-2)^{-B_h}}
 {C_\Gamma2^{B_h/2}},\qquad
 \overline a_k=\frac{C_h^{\mathrm{tilt}}}{c_\Gamma}
 k^{p_h+1}M_h(\alpha)^{k-1}.
\end{gathered}
\end{equation}
Here $c_h$ and $C_h^{\mathrm{tilt}}$ are the original TW7 and BT12
constants, as reproduced with their provider locators in \cite{NP}.
They are not free parameters chosen to make a desired conclusion true.
The original measure and its parity lift are
\begin{equation}\tag{NR4}\label{nr:measures}
 \begin{gathered}
 d\nu_0(x)=m_k(\sqrt x)x^{-1/2}\,dx,\quad d\nu_1=x\,d\nu_0,\\
 s(x)=\frac{\tanh(\alpha\sqrt x)}{\sqrt x},\quad s(0)=\alpha,
 \qquad d\rho_b=s\,d\nu_b\quad(b=0,1).
 \end{gathered}
\end{equation}
Evenness, including the Jacobian $dx=2y\,dy$ on $y>0$, gives the
exact identities for every polynomial $R$:
\begin{equation}\tag{NR5}\label{nr:lifts}
 \int_0^\infty|R(x)|^2d\rho_b(x)
 =\int_\R|y^bR(y^2)|^2s(y^2)m_k(y)\,dy.
\end{equation}
Thus both lifts preserve the actual integrals. In particular the even
untilted mass is $(\int w_h)^k$, and no mass is assigned the value one.

\section{The exact scalar minimum}

Fix either $b$ and temporarily write $\rho=\rho_b$.
Its density is positive almost everywhere and has all polynomial
moments by \eqref{nr:envelope}. Let $p_r$ be the monic degree-$r$
orthogonal polynomial for $\rho$, defined by its positive Gram system.
For $a>0$, it does not vanish at $-a$: otherwise
$q=p_r/(x+a)\in\Pcal_{r-1}$ would give
$0=\int p_rq\,d\rho=\int(x+a)q^2d\rho>0$.

\begin{lemma}[The original error minimum]
The exact scalar in RR8 is
\begin{equation}\tag{NR6}\label{nr:min}
 \lambda_r^{(b)}(a)
 =\frac{\displaystyle\int p_r(x)^2(x+a)^{-1}\,d\rho_b(x)}{p_r(-a)^2}
 =\min_{\substack{R\in\Pcal_r\\R(-a)=1}}
       \int\frac{|R(x)|^2}{x+a}\,d\rho_b(x).
\end{equation}
\end{lemma}
\begin{proof}
Put $R_*=p_r/p_r(-a)$. Every feasible difference is $(x+a)Q$ with
$Q\in\Pcal_{r-1}$. Orthogonality, valid also for complex $Q$, gives
\[
 \int\frac{|R_*+(x+a)Q|^2}{x+a}\,d\rho
 =\int\frac{|R_*|^2}{x+a}\,d\rho
    +\int(x+a)|Q|^2\,d\rho.
\]
The second term is strictly positive unless $Q=0$. This proves the
minimum and uniqueness, including $r=0$ where there is no variation.
For the matrix error, solve the polynomial Galerkin equations for
$f\in\Pcal_n$, $r\ge n+1$. Its residual $f-(x+a)q$ is orthogonal
to $\Pcal_{r-1}$ and has degree at most $r$; evaluation at $-a$ makes
it $f(-a)p_r/p_r(-a)$. Expansion of the attained square gives
\begin{equation}\tag{NR7}\label{nr:rankone}
 F_n(a)-F^-_{n,r}(a)=\lambda_r^{(b)}(a)z_n(a)z_n(a)^*,
 \quad z_n(a)=(1,-a,\ldots,(-a)^n)^T.
\end{equation}
Here $F_n(a)=\int v_n(x)v_n(x)^*/(x+a)\,d\rho_b(x)$ and
$F^-_{n,r}=C_{r,n}^*(D_r+aH_r)^{-1}C_{r,n}$, with
$H_r=[\mu_{i+j}]_{i,j<r}$, $D_r=[\mu_{i+j+1}]_{i,j<r}$,
$C_{r,n}=[\mu_{i+j}]_{i<r,j\le n}$ and
$\mu_d=\int x^d\,d\rho_b$. This identifies the error in the original
monomials, rather than only its trace or norm.
\end{proof}

\section{Gamma polynomials with their original mass}

The human source for the generating identity is
Koornwinder, Wong, Koekoek, Swarttouw and Reinhardt,
DLMF 18.23.7 \cite{DLMF18}. The precise specialization is
$p_n^\Gamma(y)=n!P_n^{(1/4)}(y/2;\pi/2)$.
We derive both the norms and the mass to fix every scale.
Euler's beta integral, in the form recorded by Askey and Roy
\cite{DLMF5}, with $t=e^u/(1+e^u)$, gives for $a>0$
\[
 |\Gamma(a+iv)|^2
 =\Gamma(2a)\int_\R e^{ivu}(2\cosh(u/2))^{-2a}\,du.
\]
Fourier inversion gives
\[
 \int_\R|\Gamma(a+iv)|^2e^{i\eta v}\,dv
 =2\pi\Gamma(2a)(2\cosh(\eta/2))^{-2a}.
\]
Both Fourier factors are integrable; for the Gamma factor this also
follows from the established exponential Gamma envelope in \cite{NP}.
Take $a=1/4$, $y=2v$ with its Jacobian, and then continue to the
convergent Laplace strip. One obtains
\begin{equation}\tag{NR8}\label{nr:transform}
 \int_\R e^{ity}\sigma(y)\,dy
 =M_\Gamma(\cosh t)^{-1/2},\qquad
 \int_\R e^{uy}\sigma(y)\,dy=M_\Gamma(\cos u)^{-1/2},\quad |u|<\pi/2,
 \qquad M_\Gamma=\sqrt{2\pi}.
\end{equation}
The branch is positive at zero. The defining generating function is
\begin{equation}\tag{NR9}\label{nr:gen}
 F(y,t)=\sum_{n\ge0}\frac{p_n^\Gamma(y)}{n!}t^n
       =(1+t^2)^{-1/4}e^{y\arctan t}.
\end{equation}
For sufficiently small real $t,u$, the cosine addition formula gives
$\int F(y,t)F(y,u)\sigma(y)dy=M_\Gamma(1-tu)^{-1/2}$.
The exponential integrability in \eqref{nr:transform} permits taking
every derivative at $t=u=0$. Therefore
\begin{equation}\tag{NR10}\label{nr:normrec}
 \begin{gathered}
 \int p_n^\Gamma p_l^\Gamma\sigma(y)\,dy
 =\mathbf1_{n=l}M_\Gamma n!(1/2)_n,\\
 p_0^\Gamma=1,\quad p_1^\Gamma=y,\quad
 p_{n+1}^\Gamma=yp_n^\Gamma-n(n-1/2)p_{n-1}^\Gamma.
 \end{gathered}
\end{equation}
The recurrence follows by differentiating \eqref{nr:gen}. It proves
parity as well as monicity.

For $j\ge1$, let $D_j=(5/4)_{j-1}$ and
$R_n(Y)=i^{-n}p_n^\Gamma(iY)$. Replacing $t$ by $-it$ in
\eqref{nr:gen} gives, at precisely $Y=2j$,
\begin{equation}\tag{NR11}\label{nr:positivegen}
 \sum_{n\ge0}\frac{R_n(2j)}{n!}t^n
 =(1+t)^{2j}(1-t^2)^{-(j+1/4)}.
\end{equation}
All coefficients are nonnegative. Retaining the constant and the
linear coefficients of the first factor gives
\begin{equation}\tag{NR12}\label{nr:coefflower}
 \frac{R_{2l}(2j)}{(2l)!}\ge\frac{(j+1/4)_l}{l!},\qquad
 \frac{R_{2l+1}(2j)}{(2l+1)!}\ge
 2j\frac{(j+1/4)_l}{l!}.
\end{equation}
For use at every integer, the two product bounds are
\begin{equation}\tag{NR13}\label{nr:products}
 \frac{n!}{(1/2)_n}\ge\sqrt{n+1},\qquad
 \frac{(5/4)_l}{l!}\ge(l+1)^{1/4},\qquad
 \frac{(j+1/4)_l}{l!}\ge\frac{(l+1)^{j-3/4}}{D_j}.
\end{equation}
The first bound starts at $n=0$; its induction step follows by squaring
from
$(n+1)^3-(n+1/2)^2(n+2)=3n/4+1/2>0$.
The second starts at $l=0$ and uses
$(1+u)^{1/4}\le1+u/4$ for $u\ge0$ at $u=1/(l+1)$.
For the third, the exact factor identity is
\[
 \frac{(j+1/4)_l}{l!}
 =\frac{(5/4)_l}{l!}\frac{(l+5/4)_{j-1}}{(5/4)_{j-1}};
\]
each numerator factor in the second ratio is at least $l+1$.

Define the two actual Gamma evaluation kernels, without combining
their different parity degrees, by
\begin{equation}\tag{NR14}\label{nr:kernels}
 K_r^{\mathrm{ev}}(2ij)=\sum_{l=0}^r
 \frac{|p_{2l}^\Gamma(2ij)|^2}{M_\Gamma(2l)!(1/2)_{2l}},\qquad
 K_r^{\mathrm{odd}}(2ij)=\sum_{l=0}^r
 \frac{|p_{2l+1}^\Gamma(2ij)|^2}{M_\Gamma(2l+1)!(1/2)_{2l+1}}.
\end{equation}
Equations \eqref{nr:coefflower}--\eqref{nr:products} bound each even
summand from below by $(l+1)^{2j-1}/(M_\Gamma D_j^2)$, because
$\sqrt{2l+1}\ge\sqrt{l+1}$. The odd summand is at least
$4j^2(l+1)^{2j-1}/(M_\Gamma D_j^2)$, since
$\sqrt{2l+2}\ge\sqrt{l+1}$. Finally comparison on each interval
$[l,l+1]$ proves
$\sum_{l=0}^r(l+1)^{2j-1}\ge(r+1)^{2j}/(2j)$.
Consequently
\begin{equation}\tag{NR15}\label{nr:kernellower}
 K_r^{\mathrm{ev}}(2ij)\ge\frac{(r+1)^{2j}}{2jM_\Gamma D_j^2},
 \qquad
 K_r^{\mathrm{odd}}(2ij)\ge\frac{2j(r+1)^{2j}}{M_\Gamma D_j^2}.
\end{equation}
The stronger odd coefficient $\sqrt2$ could be retained, but the
displayed common rate below uses precisely the incoming constant.

\section{A quantitative rate on the original native measures}

In a finite polynomial space with orthogonal basis $q_l$, evaluation
at $z$ has squared norm $K(z)=\sum_l|q_l(z)|^2/\norm{q_l}^2$.
This follows directly by Cauchy--Schwarz in the coefficients, and
equality is attained by coefficients proportional to
$\overline{q_l(z)}/\norm{q_l}^2$. Thus the minimum squared norm
with prescribed value $c$ at $z$ is $|c|^2/K(z)$.
Apply this exact fact to the even and odd spaces in
\eqref{nr:kernels}. The lifts in \eqref{nr:lifts} send $R(-a)=1$
to values $1$ and $i\sqrt a$, respectively. Since
$s(x)\le\alpha$, $x+a\ge a$ and $m_k\le\overline a_k\sigma$,
the native minimum \eqref{nr:min} satisfies
\begin{equation}\tag{NR16}\label{nr:trial}
 \lambda_r^{(0)}(a)\le
 \frac{\alpha\overline a_k}{aK_r^{\mathrm{ev}}(i\sqrt a)},\qquad
 \lambda_r^{(1)}(a)\le
 \frac{\alpha\overline a_k}{K_r^{\mathrm{odd}}(i\sqrt a)}.
\end{equation}
The $a$ in the squared odd evaluation is the reason that the
second denominator has no additional factor $a$.

\begin{theorem}[Original-pole degree rate]
For every $j\ge1$, $r\ge0$ and $b\in\{0,1\}$,
\begin{equation}\tag{NR17}\label{nr:rate}
 0\le\lambda_r^{(b)}(4j^2)
 \le\frac{\mathcal K_{k,j}}{(r+1)^{2j}},\qquad
 \mathcal K_{k,j}=
 \frac{\pi\sqrt{2\pi}}{4j}\overline a_k((5/4)_{j-1})^2.
\end{equation}
\end{theorem}
\begin{proof}
Substitute $a=4j^2$ and \eqref{nr:kernellower} into
\eqref{nr:trial}. Both prefactors are
$\alpha\overline a_k M_\Gamma D_j^2/(2j)$, which is the displayed
$\mathcal K_{k,j}$. Every trial integral before this comparison
is still an integral against its original native measure.
\end{proof}

\section{An explicit sufficient degree for the finite inner test}

Let $N\ge0$ be the original total polynomial degree,
$n_0=\lfloor N/2\rfloor$, and $n_1=\lfloor(N-1)/2\rfloor$.
Omit the odd block when $N=0$. Put
\[
 H_{\rho_b}=[\textstyle\int x^{i+l}d\rho_b]_{i,l=0}^{n_b},\qquad
 H_{\mathrm{base}}=\frac2\pi
       \operatorname{diag}(H_{\rho_0},H_{\rho_1})>0.
\]
Write $z_{b,j}$ for $z_{n_b}(a_j)$ padded into parity block $b$,
$a_j=4j^2$, and
$\chi_{b,j}=z_{b,j}^*H_{\mathrm{base}}^{-1}z_{b,j}$.
For a finite retained pole range $1\le j\le J$ choose any positive
$\varepsilon_{b,j}$ with sum at most $\delta_{\mathrm{inner}}$.
These are requested accuracy allocations, not assignments to the
programme's unresolved arithmetic kernel and complementary volumes.

The original finite-moment Radau width is
\[
 \gamma_r(a)=\frac{\kappa_r}
 {p_r(-a)^2(1-\kappa_re^*(D_r+aH_r)^{-1}e)},\qquad
 \kappa_r=(e^*D_r^{-1}e)^{-1}.
\]
Here $e$ is the final monomial vector, and RR7--9 prove
$0\le\lambda_r\le\gamma_r$ and a positive denominator.
Set
\begin{equation}\tag{NR18}\label{nr:scalarcap}
 \widehat\gamma_{r,j}
 =\min\{\gamma_r(a_j),\mathcal K_{k,j}(r+1)^{-2j}\}.
\end{equation}
This minimum is of two scalar bounds along the same proved direction
$zz^*$. Since $B_n(a)=H_{\rho_b}-aF_n(a)$, its new lower matrix is
\[
 B^{\mathrm{new},-}_{n,r}(a_j)
 =H_{\rho_b}-a_jF^-_{n,r}(a_j)
       -a_j\widehat\gamma_{r,j}z_n(a_j)z_n(a_j)^*.
\]
It is below $B_n(a_j)$ by \eqref{nr:rankone} and
\eqref{nr:scalarcap}. It is above the old Radau lower matrix
$H_{\rho_b}-a_jF^+_{n,r}(a_j)\succeq0$, because
$\widehat\gamma\le\gamma$ and $F^+-F^-=\gamma zz^*$.
Thus it is positive semidefinite.

The explicit degree is
\begin{equation}\tag{NR19}\label{nr:degree}
 r_{b,j}=\max\left\{n_b+1,
 \left\lceil\left(
 \frac{4a_j\mathcal K_{k,j}\chi_{b,j}}
 {\pi\varepsilon_{b,j}}\right)^{1/(2j)}\right\rceil-1\right\}.
\end{equation}
For any positive $H$, Cauchy--Schwarz applied to
$H^{-1/2}z$ and $H^{1/2}u$ proves
$zz^*\preceq(z^*H^{-1}z)H$. Applying this with
$H=H_{\mathrm{base}}$, and then \eqref{nr:degree}, proves
\begin{equation}\tag{NR20}\label{nr:inner}
 \frac4\pi\sum_{b,j}a_j\widehat\gamma_{r_{b,j},j}
 z_{b,j}z_{b,j}^*
 \preceq\delta_{\mathrm{inner}}H_{\mathrm{base}}
 \preceq\delta_{\mathrm{inner}}G^-.
\end{equation}
Here $G^-$ is either the old or new finite source lower Gram: both
are $H_{\mathrm{base}}$ plus positive semidefinite pole summands.
The outer pole tail is unchanged. No assertion about the sign of a
complex arithmetic current is needed for this inner estimate.

\section{A proved bound on the original inverse evaluation}

For every $y\in\R$,
\begin{equation}\tag{NR21}\label{nr:slow}
 s(y^2)\ge\tanh(\alpha)(1+y^2)^{-1/2}
 \ge\frac{\tanh(\alpha)}{\sqrt2}(1+y^2)^{-1/2}.
\end{equation}
Indeed $\tanh(\alpha t)/t$ decreases for $t>0$:
$\tanh v-v\operatorname{sech}^2v$ has derivative
$2v\operatorname{sech}^2v\tanh v\ge0$ and vanishes at zero.
For $0\le|y|\le1$ this gives $s(y^2)\ge\tanh\alpha$.
For $|y|\ge1$, use $\tanh(\alpha|y|)\ge\tanh\alpha$ and
$|y|\le\sqrt{1+y^2}$. The weaker second constant is exactly the
constant in the incoming continuation.

For $P\in\Pcal_N$, expand in the orthonormal Gamma basis.
The recurrence \eqref{nr:normrec} splits multiplication by $y$
into an upward and a downward weighted shift. Their respective norms
on this coefficient space are at most
$\sqrt{(N+1)(N+1/2)}$ and $\sqrt{N(N-1/2)}$, taking the latter
as zero when $N=0$. The triangle inequality therefore gives
\begin{equation}\tag{NR22}\label{nr:mult}
 \int y^2|P(y)|^2\sigma(y)\,dy
 \le4(N+1)^2\int|P(y)|^2\sigma(y)\,dy.
\end{equation}
This is an estimate of the full multiplication image in
$\Pcal_{N+1}$, so no top coefficient is discarded.

Put $L=M+1/2>0$. For nonzero $P$, apply the tangent-line inequality
for the convex function $f(t)=(1+t)^{-L}$ at
$t_0=\int y^2|P|^2\sigma/\int|P|^2\sigma$ and integrate against
$|P|^2\sigma$. The linear term has integral zero. Using
\eqref{nr:mult} and the decreasing property of $f$ gives
\begin{equation}\tag{NR23}\label{nr:jensen}
 \int|P(y)|^2(1+y^2)^{-L}\sigma(y)\,dy
 \ge[1+4(N+1)^2]^{-L}\int|P(y)|^2\sigma(y)\,dy.
\end{equation}
Both sides retain the original mass. With
\begin{equation}\tag{NR24}\label{nr:lowerconstant}
 c_{k,N}=\frac{\tanh(\pi/2)}{\sqrt2}\underline a_k
        [1+4(N+1)^2]^{-(M+1/2)},
\end{equation}
equations \eqref{nr:envelope}, \eqref{nr:slow} and
\eqref{nr:jensen} prove, for either original parity block,
\begin{equation}\tag{NR25}\label{nr:gramlower}
 \int|R(x)|^2d\rho_b(x)
 \ge c_{k,N}\int_\R|y^bR(y^2)|^2\sigma(y)\,dy,
 \qquad \deg R\le n_b,\quad 2n_b+b\le N.
\end{equation}
Inverting this positive Gram comparison and applying the exact
evaluation-kernel identity proves
\begin{equation}\tag{NR26}\label{nr:chi}
 \boxed{\chi_{b,j}\le
 \frac{\pi}{2c_{k,N}a_j^b}K_{n_b}^{(b)}(2ij),\qquad
 K^{(0)}=K^{\mathrm{ev}},\quad K^{(1)}=K^{\mathrm{odd}}.}
\end{equation}
The factor $a_j^{-b}$ is forced by the exact value of the parity
lift at $2ij$, not by a convention on the measure. The condition
$2n_b+b\le N$ is essential to the multiplication estimate's index;
it is automatically satisfied by the original cutoff definitions.
Using the first, stronger inequality in \eqref{nr:slow} would
replace $c_{k,N}$ by $\sqrt2c_{k,N}$ and improve this upper bound
by $\sqrt2$. The incoming bound is correct as written with its
smaller constant.

Every number on the right of \eqref{nr:chi} is a finite expression
in the original envelope constants and the Gamma recurrence. It
can replace $\chi_{b,j}$ in \eqref{nr:degree} without weakening the
guarantee. This proves a finite sufficient degree; it does not assign
numerical values to the arithmetic envelope constants, evaluate the
native moments, or establish the unresolved projected-current sign.

\begin{thebibliography}{9}
\bibitem{ZDS} J.~Zimmerling, V.~Druskin and V.~Simoncini,
\emph{Monotonicity, bounds and acceleration of Block Gauss and
Gauss--Radau quadrature for computing $B^T\phi(A)B$},
arXiv:2407.21505v3, original author LaTeX, Sections 4--5.1.
\url{https://arxiv.org/abs/2407.21505v3}.
The title is transcribed from the version-3 author source.
\bibitem{DLMF18} T.~H.~Koornwinder, R.~Wong, R.~Koekoek,
R.~F.~Swarttouw and W.~P.~Reinhardt, \emph{Orthogonal Polynomials},
NIST Digital Library of Mathematical Functions, Chapter 18,
18.23.7, version 1.2.8 (15 September 2026).
\url{https://dlmf.nist.gov/18.23.E7}.
The equation's original TeX was read. Its source notes credit
M.~E.~H.~Ismail (2009), (5.9.3); that book was not separately read
for this calculation.
\bibitem{DLMF5} R.~A.~Askey and R.~Roy, \emph{Gamma Function},
NIST Digital Library of Mathematical Functions, Chapter 5,
Euler's beta integral 5.12.1, version 1.2.8 (15 September 2026).
\url{https://dlmf.nist.gov/5.12.E1}. Original equation TeX read.
\bibitem{RR} Split-Zero programme,
\emph{One scalar resolvent uncertainty in the original polynomial observation},
20 September 2026, RR1--16. Local full proof:
\texttt{NATIVE\_RESOLVENT\_RANK\_ONE.tex}.
The original measures, rank-one direction and Radau scalar are from
this earlier programme derivation; NR6--7 supply an independent
variational proof of the portion used here.
\bibitem{NP} Split-Zero programme,
\emph{Explicit native moment precision, inverse margins, and the original
signed receiver}, 20 September 2026, first section, equations
\texttt{eq:objects}, \texttt{eq:envelope}, \texttt{eq:measures}.
Local full proof: \texttt{NATIVE\_MOMENT\_PRECISION.tex}.
These equations retain the original CAI16--17 and HG21
arithmetic envelope and its TW7/BT12 constants.
\end{thebibliography}
\end{document}
